\cotefile{136}
\scanpage{48}
% @01–02, upper right, boxed heading
\fbox{\shortstack{Cohomologie des ens. $1/2$ simpliciaux\\et \unsure{résolutions comme $1/2$ simpliciales|résolution comme $1/2$ simpliciale|résolutions canon. $1/2$ simpliciales} ?}}

% @03, opening
Soit $k$ un \deleted{groupe abélien}\inserted{anneau}, $M$ un $k$-module,
% @04, functor
\[
 C^*\colon\widehat{\Delta}^{\,\circ}\longrightarrow
 \text{complexes de cochaines de $k$-modules}
\]
% @05, hypotheses
un foncteur satisfaisant les \unsure{suppos.}

% @06, condition a
a) $\forall\ X_*\in\operatorname{Ob}\widehat{\Delta},\quad H^*(X_*,M)\simeq H^*(C^*(X_*))$ (isom fonctoriel en $M$ -- au moins pour $H^0$)

% @07, condition b
b) $\forall n\geq 0,\quad X_*\longmapsto C^n(X_*)$ transforme les \unsure{$\varinjlim$ en $\varprojlim$}.

% @08, representability
Alors b) signifie que les $C^n$ sont représentables
% @09, representation
\[
 C^n(X_*)=\operatorname{Hom}(X_*,C_*^n).
\]
% @10–14, resolution and augmentation
a) implique, en faisant $X_*=\Delta^m$, que les $C_{\unsure{m|k|\bullet}}^*$ \unsure{sont des} \uline{résolutions} de $\uline{M}$, alors que l'augmentation $M\to C_m^*$ \unsure{soit} fonctorielle en $m$, i.e. $C_*^*$ est une résolution de $M_*=\unsure{S_*|\underline{k}}\otimes_k M\simeq M_{\unsure{s.s.}}$ \unsure{ou $M$}, ce qui explique bien que l'on a, \unsure{aux conditions de} fonctorialité \unsure{près}, $H^*(C^*(X_*))=H^*(X_*,M)$,
\unsure{des suites spectrales}. \illegible{2} pour $X_*$ \unsure{de la forme} $\Delta^m$.
% @14a, left margin
\margin{\unsure{Cela l'intègre au programme ?}}

% @15–18, converse across the fold
\deleted{En prenant} $X_*$ \unsure{``arbitraire''} on trouve que \unsure{$H^k(X_*,M)$} \illegible{3} $\unsure{H^0(C^*(X_*))}=M$ \unsure{pour $H$}. \unsure{Il est a priori} \unsure{faussement|fraîchement} fort que $C_*^\bullet$ \unsure{soit} une résolution de $M_*$) i.e. \unsure{tout se trouve} $X_*\to Z_*^i$.
% @19, two lifting diagrams
\diagram{\begin{tikzcd}[ampersand replacement=\&,column sep=small,row sep=small]
 C_*^{i-1}\arrow[r]\&Z_*^i\&[-1em]\text{(\,pour}\&C_*^0\arrow[r]\&Z_*^1\&\text{)}\\
 \& X_*\arrow[u]\&\&\&X_*\arrow[u]\&
\end{tikzcd}}
% @20–23, lifting and flasqueness
se relève en $C_*^{i-1}$ \unsure{--}. On a envie de prouver que a) pour des quelconques signifie que les $C_*^i$ sont flasques, (ce qui serait \unsure{évidemment} suffisant pour qu'on ait a)).
% @23–25, cancelled lines
\deleted{\unsure{ce qu'il y a là} \illegible{3} \unsure{ce qui} \unsure{est démontré} que \illegible{2} fournissant \unsure{alors} a) \unsure{en résulte}. \inserted{\unsure{(cf. question dessus)}} \unsure{appliquer a)} \illegible{1}. On aura}
% @26, cohomology row
\[
 H^1(X_*,C_*^0)\to H^0(X_*,Z_*^1)\to H^1(X_*,M)\to H^1(X_*,C_*^0)\xrightarrow{\unsure{\text{inj}}} H^1(X_*,Z_*^1).
\]
% @27–29, paragraph
Mais c'est \unsure{faux} déjà pour $M=0$. Il suffit, pour qu'on ait \unsure{b)|a)}, que la suite exacte \unsure{longue} des $C_*^{\unsure{0|n}}\to Z_*^{\unsure{1|i}}\to\cdots$ soit ``splittée'' \unsure{trivialement|bêtement} : \unsure{homotopie} !

% @30–31, observation
\uline{NB} Si on a $C_*^\bullet$ une résolution flasque de $M_*$, les $Z_*^n(C_*^\bullet)$ sont des rétractes des $K_*(M,n)$ \ldots
% archive sheet number: 47

\scanpage{49}
% typed verso
% @top right, pagination
\mspage{6}
% @01–02, opening
Continuing in this way we construct the minimal model
% @03, first display
\[
 \Lambda(x_1,x_2,\ldots\mathbin{:}d)\underset{\rho}{\to}\mathcal{E}
\]
% @04, continuation
satisfying properties i) and ii) above.

% @05–06, theorem heading
\uline{Theorem B}: Let $X$ be a triangulated space which is simply connected. If
% @07, second display
\[
 \Lambda(x_1,x_2,\ldots)\underset{\rho}{\to}\mathcal{E}
\]
% @08, hypothesis
is a minimal model of the de Rham algebra then

% @09–10, conclusion i
i) the rank of $\pi_i X$ is the number of variables in degree $i$.

% @11–12, conclusion ii
ii) the rank of the Hurewicz homomorphism $\pi_i X\to H_i X$ is the number of $d$-closed generators in degree $i$.

% @13–15, conclusion iii
iii) the Whitehead products on homotopy are described by the quadratic terms of the $d$-formulae of generators.

% @16–23, proof discussion
Theorem B is proved by a not so standard Postnikov argument using Theorem A and the commutativity of the wedge product multiplication in $\mathcal{E}$. The argument is based on the beautifully simple Guy Hirsh method of computing the cohomology of a principal $G$-bundle over a manifold using differential forms (see appendix H). The Hirsch method was suggested by Rene Thom in the study of Postnikov systems in the Cartan Seminar 1954.

% @24–25, final paragraph
The proof of theorem B shows the model can be used to directly construct the \uline{rational} Postnikov system
% archive sheet number: 48

\scanpage{50}
% @01–02, opening example
\uline{Ex 1} \unsure{N°|Cf.} \unsure{ici}, p. \unsure{35}, \unsure{en général un foncteur fibre|on avait un foncteur fibre} $\unsure{v|u}_n^*(K)=K_n$
% @03, chain of identifications
\[
 M\mapsto C_{\unsure{v}}^n(K_*,M)\overset{\mathrm{dfn}}{=}
 \operatorname{Hom}_{\unsure{\mathrm{Ens}|\widehat{\Delta}}}(K_n,M)
 \simeq\operatorname{Hom}_{\mathrm{Ens}}(\unsure{v_n^*}(K_*),M)
 \simeq\operatorname{Hom}_{\widehat{\Delta}}(K_*,\unsure{v_{n*}}(M)).
\]
% @04–08, representing object
est représentable par \unsure{$v_{n*}M$}. Plus \unsure{précisément}, les \unsure{$v_{n*}(M)$} forment un $k$-Module \unsure{cosimplicial} de $\widehat{\Delta}$, d'où (\unsure{ce qui a la rigueur s'exprime|vu que à la source} \unsure{avec $\varphi$ et fon}) un complexe de cochaines de $k$-Modules, qui est $C_{*M}^{*}$. D'où
% @09, cochain functor
\[
 C^n(K_*,M)\simeq\operatorname{Hom}(K_*,C_{*M}^n).
\]
% @10–15, resolution
\unsure{Comme on doit|Considérons donc}, de $k$-Modules cosimpliciaux \unsure{en} des complexes de \unsure{cochaines de $k$-modules} (fonct. en \unsure{$K_*$|$K_\bullet$}) de sorte que pour $K_*$ \unsure{de base}, on trouve une \uline{résolution} de $M$, i.e. $C_{*M}^\bullet$ est une résolution de $M_*$. De plus les $C_{*M}^n=\unsure{v_{n*}}(M)$ sont flasques \unsure{car les faisceaux $M$ sur Pt le sont} !
% @15a, small right-hand note
\margin{NB \unsure{$C_{mM}^n=v_{n*}(M)=\Delta^m\longmapsto\operatorname{Hom}(\Delta_n^m,M)$} $\unsure{=M^{\operatorname{Hom}(\Delta_n,\Delta_m)}}\simeq{k^{\operatorname{Hom}(\Delta_n,\Delta_m)}\otimes_k M}$.}

% @16–18, second example
\uline{Ex 2} Considérons le sous-complexe de cochaines de $k$-Modules de $C_{*M}^\bullet$, formé des cochaines non dégénérées :
% @19–22, normalization
\[
\begin{aligned}
 C_{*M}^{\bullet\prime}&=(C_{*M}^{n\prime})_{n\geq0}\text{ où}\\
 C_{*M}^{n\prime}&=\bigcap_{\substack{\Delta_n\xrightarrow{\sigma}\Delta_m\\ \sigma\ \unsure{\text{épi}},\ m<n}}\operatorname{Ker}(\sigma_*\colon C_{*M}^{n}\longrightarrow C_{*M}^{m})\\
 &=\bigcap_{\substack{\unsure{h_i}\colon\Delta_n\to\Delta_{n-1}\\ \unsure{\text{épi}|i,j}}}\operatorname{Ker}(\unsure{h_i}\colon C_{*M}^n\longrightarrow C_{*M}^{n-1})\\
 &=\bigcap_{\substack{\Delta_n\xrightarrow{\sigma}\Delta_m\\m<n,\ \unsure{\text{tout }\sigma}}}\operatorname{Ker}(\sigma_*\colon C_{*M}^n\longrightarrow C_{*M}^m).
\end{aligned}
\]
% @23–25, final paragraph
On sait que cette opération fournit, par \unsure{rétraction} sur ce \unsure{facteur direct} qui \unsure{résulte|se calcule} \unsure{ -- cela vaut dans un cadre plus général} qui nous \unsure{assure} $C^{\bullet\prime}(X_*)$ \unsure{avec} $C^\bullet(X_*)$,
% archive sheet number: 49

\scanpage{51}
% typed verso
% @top right, pagination
\mspage{5}
% @01–04, continuation
algebra of $X$ (in any sense), then one finds the second homotopy \unsure{grou|group} $\pi_2 X$ has rank $n$, $\pi_3 X$ has rank $\ell+m$, the Hurewicz map $\pi_3 X\to H_3 X$ has rank $\ell$, and the symmetric pairing given by Whitehead product
% @05, bracket
\[
 \pi_2 X\handwritten{\otimes}\pi_2 X\mathrel{\overset{[,\!]}{\handwritten{\longrightarrow}}}\pi_3 X
\]
% @06–07, tensor
is described over the rationals by the $i,j$ symmetric tensor $(a_{ij}^k)$.)

% @08, induction
Now suppose by induction that
% @09, model and map
\[
 \Lambda(x_1,\ldots,x_n;y_1,y_2,\ldots;z_1,z_2,\ldots:d)
 \mathrel{\overset{\rho_n}{\handwritten{\longrightarrow}}}\handwritten{\mathscr{E}}
\]
% @10–13, induction hypothesis
has been constructed so that $d$ of any generator is a polynomical in previous generators while $\rho_n$ is an isomorphism on cohomology up through dimension $n$ and injective in dimension $n+1$.

% @14–25, adjoining variables
Then as in step 2 we adjoin variables in dimension $n+1$ to make $\rho$ onto in dimension $n+1$ and injective in dimension $n+2$. In this general case we add variables $\omega_1,\ldots,\omega_d,\omega_{d+1},\ldots,\omega_r$ and define $d$ and $\rho=\rho_{n+1}$ as above. Now $d\omega_1,\ldots,d\omega_d$ are zero and $\rho\omega_1,\ldots,\rho\omega_d$ give generators of the cokernel of $\rho_n$ on $H^{n+1}$, $d\omega_{d+1},\ldots,d\omega_r$ are polynomials in $x_1,x_2,\ldots;y_1,y_2,\ldots;z_1,z_2,\ldots$ which generate the kernel of $\rho_n$ on $H^{n+2}(\Lambda(x_1,x_2,\ldots;y_1,y_2,\ldots;z_1,z_2,\ldots:d))$ and $\rho\omega_{d+1},\ldots,\rho\omega_r$ are elements in $\mathcal{E}$ whose differentials are these same polynomials in
% @26, final display
\[
 \rho x_1,\rho x_2,\ldots;\rho y_1,\rho y_2,\ldots;\rho z_1,\rho z_2,\ldots\ .
\]
% archive sheet number: 50

\scanpage{52}
% @01–04, continuation
\unsure{vu|et} que l'opération $^{!}$ de prendre les \unsure{sous-complexes} des cochaines non dégénérées commute à l'op. $\operatorname{Hom}(X_*,-)$. Donc $C^{\bullet!}(X_*)=\operatorname{Hom}(X_*,C_{*M}^{\bullet!})$ satisfait aux cond. a) et b). \unsure{Il est remarquable que|Il est connu que} $C_{*M}^{n!}$ \unsure{soit flasque} : en effet, on explicite aisément :
% @06–07, normalized cochains
\[
\begin{aligned}
 (C_m^{n!})_M&\simeq M^{\operatorname{Hom\,mon}(\Delta_n,\Delta_m)}
 \simeq k^{\operatorname{Hom}^{\mathrm{mono}}(\Delta_n,\Delta_m)}\otimes_k M\\
 &\simeq(\bigwedge\nolimits^{n+1} k^{\Delta_m})\otimes_k M
\end{aligned}
\]
% @08, functoriality
et on vérifie que \unsure{l'on a fonctoriellement|l'on a fonctorialité} sur $\Delta^\circ$
% @09, simplicial form
\[
 C_{*M}^{n!}\simeq\bigwedge\nolimits^{n+1}\Phi_{*k}\otimes_k M_*,
\]
% @10–11, definition
où $\unsure{\Phi_{*k}|\underline{\Phi}_{*k}}=(\Delta_m\longmapsto k^{\Delta_m})$. [On vérifie \unsure{des} \unsure{plus|faits} ceci : soit
% @12, map T
\[
 T\colon k_*\longrightarrow\Phi_{*k}=C_{*k}^{0!}
\]
% @13–18, bracketed discussion
l'élém. unité (provenant de $\Phi_{\unsure{01}}\simeq k$). Alors l'\unsure{opération différent.|opérateur différentiel} dans $C_{*M}^{\bullet!}\simeq\bigwedge\Phi_{*k}\otimes_k M$ est le produit par $T\wedge$ (c'est un fait de \unsure{ch} \illegible{2} pour des calculs de \unsure{dégénérescence})]. Or, $\Phi_{*k}$ est flasque et a la même propriété, de \unsure{tautologie} ?, les \unsure{somm} $\bigwedge^i\Phi_{*k}\otimes_k M$. Le complexe $\bigwedge(T)\otimes_k M$, pour $n>0$, est \deleted{\unsure{ex}} exact, \unsure{etc}.

% @20–24, final paragraph
\unsure{Aussi bien} au lieu de vérifier la flasquité des $C_{*M}^{n!}$ : c'est, \deleted{\unsure{en}} \unsure{naturel de prouver directement que} \illegible{2} des \unsure{rétractés} des $C_{*M}^n$ (\unsure{savoir|Savoir} que la \unsure{rétraction} \unsure{engendrée par les $\delta_i^*$} \inserted{\unsure{morph.}} : $C_{*M}^m\to C_{*M}^{\unsure{n}}$ pour $\unsure{i,n}$ sont \unsure{transformées}).
% archive sheet number: 51

\scanpage{53}
% typed verso
% @top right, pagination
\mspage{4}
% @01–03, assumptions
Suppose always that $\mathcal{E}$ is homologically connected and for the moment that the first Betti number of $\mathcal{E}$ is zero.

% @04–08, first step
\uline{Step} \uline{1}: Construct $\Lambda(x_1,\ldots,x_n)\mathrel{\overset{\rho_1}{\handwritten{\longrightarrow}}}\mathcal{E}$ where $x_1,\mathinner{..},x_n$ are in dimension 2 and $\rho_1x_1,\ldots,\rho_1x_n$ generate $H^2\mathcal{E}$. Define $d$ in $\Lambda(x_1,\ldots,x_n)$ by $dx_i=0$, $i=1,\ldots,n$. Note that $\rho_1$ is an isomorphism on cohomology in dimension 2 and injective in dimension 3.

% @09–11, second step
\uline{Step} \uline{2}: Adjoin three-dimensional variables to $\Lambda(x_1,\ldots,x_n)$ to make $\rho_1$ surjective on cohomology in dimension 3 and injective on cohomology in dimension 4.

% @12, construction
More explicitly, we form
% @13, list
\[
 \Lambda(x_i;y_j),\ \rho_2,\ \text{and }d.
\]
% @14, definition
$d$ is defined by
% @15–16, differentials
\[
\begin{aligned}
 dy_1&=0,\ldots,dy_\ell=0\\
 dy_{\ell+k}&=\sum a_{ij}^k\,x_i\wedge x_j\qquad k=1,\ldots,m
\end{aligned}
\]
% @17–18, typed slip
\inserted{where $\ell$ is the rank of $H^3\handwritten{\mathscr{E}}$ and $\{\sum a_{ij}^k x_i\wedge x_j\}$ is a basis for kernel $\rho_1$ on 4-dimensional cohomology.}

% @19–20, choosing images
$\rho_2$ is defined by choosing $\rho_2y_1,\ldots,\rho_2y_\ell$ to generate $H^3\mathcal{E}$ and $\rho_2(y_{\ell+k})$, $k=1,\ldots,m$, to satisfy
% @21, differential identity
\[
 d(\rho_2(y_{\ell+k}))=\sum_{i,j}a_{ij}^k\,\rho_1(x_i)\wedge\rho_1(x_j).
\]
% @22, conclusion
(Now $\rho_2$ is an isomorphism on $H^3$ and injective on $H^4$.

% @23–25, Postnikov argument
At this point one can show by a standard Postnikov argument how this construction relates to the homotopy theory of a simply connected space $X$. If $\mathcal{E}$ is the de Rham
% archive sheet number: 52

\scanpage{54}
% @01, underlined heading
\uline{Intégration} \uline{des Algèbres $1/2$ simpliciales}

% @02–06, graded object
Soit $A_*$ un groupe gradué différentiel $1/2$ simplicial i.e. un objet $1/2$ simplicial dans la catégorie des groupes gradués différentiels (: op. de dérivation de degré $+1$) \deleted{On suppose} Donc $\forall n\geq0$ $A_n=A_n^*=\sum_{p\in\mathbb{Z}}A_n^p$, avec op. différentiel $d_n\colon A_n^{\unsure{\bullet|*|i}}\to A_n^{\unsure{\bullet|*}}$, $A_n^p\to A_n^{p+1}$. On suppose

% @07–11, hypothesis
(1) $\forall n$, on a $A_n^p=0$ si $p\notin[0,n]$, et \deleted{\unsure{$l'$|$k$}} $A_n^*$ muni de l'augmentation \inserted{$\varepsilon_n\colon$} $A_0\to A_n^*$ [\uline{NB} associée à l'unique $\sigma\colon\Delta_n\to\Delta_0$ ; $A_0$ est considéré comme groupe doué de la graduation nulle et d'op. diff trivial] est une \uline{résolution} de $A_0$.

% @12–13, boundary
Rappelons qu'on a sur le système des $A_n$ un op. différentiel (de degré $-1$)
% @14–15, two defining lines
\[
\begin{aligned}
 \partial&\colon A_n^*\longrightarrow A_{n-1}^*\\
 \partial(\omega)&=\sum_{0\leq i\leq n}\deleted{\unsure{d}}\partial_{n,i}^*(\omega)(-1)^i
\end{aligned}
\]
% @16–17, faces
(où les $\partial_{\unsure{n,i|n-i}}$ \unsure{$1\leq i\leq n$|$0\leq i\leq n$}) sont les $n+1$ applications strictmt croissantes $\Delta_{n-1}\to\Delta_n$. On sait que $\partial(\partial\omega)=0$

% @18–20, proposition
\uline{Proposition} 1 Il existe un et un seul système d'hom de groupes
% @19, maps
\[
 A_n^*\xrightarrow{k_n}A_0
\]
% @20, conditions
satisfaisant aux conditions suivantes
% @20a, left margin
% slanted
\margin{\uline{NB} le système des $k_n$ \uline{ne commute pas} en gén. aux op. $1/2$ simpliciales !}
% @21–23, conditions
\[
\begin{aligned}
 &\text{a) (normalisation)}& k_0&=\operatorname{id}_{A_0}\\
 &\text{b) (Condition de degré)}& k_{\unsure{n|m}}(A_n^p)&=0\quad\text{si }p\ne n\\
 &\text{c) (Stokes)}& k_n(d\omega)&=k_{n-1}(\partial\omega)\quad:\ \forall\omega\in A_n^\bullet.
\end{aligned}
\]
% @24–25, reduction
En vertu de b), la condition c) équivaut à la $=$ condition, avec $\omega\in A_n^{n-1}$. On montre existence et unicité \deleted{par} \inserted{\unsure{degré $\leq N$ par|de $k_n$, $n\leq N$, par|degré $\in\mathbb N$ par}} récurrence sur $N$, en notant que pour $N=0$ c'est trivial. Si $N>0$, \deleted{\unsure{si c'est|$\in A$}} il reste à démontrer ex. et un. de $k_N$ i.e. des $k_N(\overline\omega^N)$ pour $\deleted{\unsure{N}}\overline\omega^N\in A_N^N$.
% @28–30, exactness and Stokes
Or $d\overline\omega^N=0$ \deleted{\unsure{puisque}} car $A_N^{N+1}=0$, \inserted{(1)} donc $\overline\omega^N=d\omega^{N-1}$ (1). [En fait il suffit ici \deleted{\unsure{ici|i.e.|ier}} $H^N(A_N^*)=0$ pour $N\geq1$] on le \uline{définit} par c) i.e.
% @31, last equation
\[
 k_n(d\omega^{N-1})=k_{n-1}(\partial\omega)
\]
% archive sheet number: 53

\scanpage{55}
% typed verso
% @top centre, pagination
\mspage{- 130 -}
% @top right, running head
\hfill IX

% @01–02, opening
Elle implique bien que si $\varphi^t_{\xi_{\unsure{o|0}}}$ est symétrique et positive, il en est de même de $\varphi^t_\xi$.

% @03–11, nondegeneracy
Enfin, le fait que $\varphi^t_{\xi_o}$ soit non dégénérée équivaut évidemment à l'assertion que l'accouplement (10.4.1) est séparant (compte tenu que $\xi_{ot}:M\otimes_{\mathbb{Z}}\mathbb{Q}\handwritten{\longrightarrow}M'\otimes_{\mathbb{Z}}\mathbb{Q}$ est un isomorphisme, $\widetilde{\xi}_o$ étant une isogénie), et est donc également indépendant de la polarisation choisie. On peut d'ailleurs noter que pour vérifier que (10.4.1) est séparant, il suffit de vérifier la même assertion pour l'accouplement de monodromie $u_\ell$ (9.1.2) pour \uline{un} nombre premier $\ell$. Prenant $\ell\ne p$, on peut utiliser la description 9.2 de $u_\ell$ ; le fait que $u_\ell$ soit séparant à gauche (donc des deux côtés, pour une raison de rangs) est alors trivial.

% @12–13, section title
11. \uline{Composantes connexes du modèle de Néron} : \uline{relation de dualité}, \uline{comportement assymptotique}\unsure{.}

% @14, opening of section
11.0. Le but du présent paragraphe est l'étude du $k$-groupe fini étale
% @15, component group
\[
 \Phi_o=A_o/A_o^o\tag*{(11.0.1)}
\]
% @16, introducing the primary component
déjà considéré dans 1. Pour tout nombre premier $\ell$, nous désignons par
% @17, primary component
\[
 \Phi_o(\ell)
\]
% @18–19, decomposition
la composante $\ell$-primaire de ce dernier, de sorte que la connaissance de $\Phi_o$ équivaut à celle des $\Phi_o(\unsure{1|l})$, grâce à la formule
% @20, decomposition formula
\[
 \Phi_o=\handwritten{\coprod}_{\ell}\Phi_o(\ell)\quad .\tag*{(11.0.2)}
\]
% @lower strip, pencil foliation
\margin{\handwritten{448}}
% @lower centre, printed label
\mspage{448}
% archive sheet number: 54

\scanpage{56}
% @01–02, independence
il faut simplement voir que celà ne dépend pas du choix de $\omega^{N-1}$ i.e. que \deleted{$d\omega^{N-1}=0$}
% @03, implication
\[
 d\omega^{N-1}=0\Longrightarrow k_{N-1}(\partial\omega^{N-1})=0
\]
% @04, cases
Distinguons deux cas :

% @05–06, first case
a) $N\geq2$ i.e. $N-1\geq1$, alors $H^{N-1}(A_N^*]=0$ par (1), donc $d\omega^{N-1}=0\Rightarrow\omega^{N-1}=d\alpha^{N-2}$, et on aura
% @07, substitution
\[
 k_{N-1}(\partial\omega^{N-1})=k_{N-1}(\partial d\alpha^{N-2}]\unsure{=|\Leftarrow}
\]
% @08, commutation
Or $\deleted{\partial\alpha^{N-}}\ \partial d\alpha^{N-2}=d\partial\alpha^{N-2}$ donc \unsure{on trouve}
% @09, Stokes
\[
 k_{N-1}(\partial\omega^{N-1})=k_{N-1}(d\partial\alpha^{\unsure{N-1|N-2}})
 \overset{\text{Stokes}}{=}k_{\unsure{N-1|N-2}}(\underbrace{\partial\partial\alpha^{N-2}}_0)=0
\]
% @10, end of case
cqfd

% @11–12, second case
b) $N=1$ i.e. $N-1=0$. Alors en vertu de (1) \deleted{\unsure{on|on a}} $d\omega^{\deleted{\unsure{n}}0}=0$ implique $\omega^{\deleted{\unsure{n}}0}=\varepsilon_1(\alpha)$ ($\alpha\in A_0$) donc
% @13, augmentation
\[
 k_{N-1}(\partial\omega^0)=k_{N-1}(\underbrace{\partial\varepsilon_1(\alpha)}_{=0})=0
 \quad\deleted{\text{or }\partial\varepsilon_1(\alpha)\unsure{=0}}\quad\text{o.k.}
\]
% @14–15, first remark
\uline{Remarques}. 1) Dans \inserted{l'hyp.} (1), en plus de la condition des degrés, on a utilisé
% @16–17, assumptions used
\[
 \left\{\begin{aligned}
 H^N(A_N^*)&=0\quad\text{si }N\geq1\\
 H^{N-1}(A_N^*)&=\begin{cases}0&\text{si }N\geq2\\\varepsilon_1(A_0)&\text{si }N=1\end{cases}
 \end{aligned}\right.
\]
% @18–22, second remark
2) On peut donner une version ``abstraite'' en termes d'un bicomplexe augmenté avec des conditions de degré qu'il faut (et $d\partial=\partial d$ \unsure{!}) \deleted{$\partial\varepsilon_1=0$}) dans une cat. abélienne [\unsure{celà|ceci|celle} peut être $(\mathrm{Ab})^\circ$ !] Conditions
% @18–24, left-hand double complex
\diagram{\begin{tikzcd}[ampersand replacement=\&,column sep=small,row sep=small]
 \&\&\&\&0\&\\
 \&\&\&0\&A_3^3\arrow[u]\&{}\arrow[l]\\
 \&\&0\& A_2^2\arrow[u]\arrow[l]\& A_3^2\arrow[u,"d"']\arrow[l,"\partial"']\&{}\arrow[l]\\
 \&0\& A_1^1\arrow[u]\arrow[l]\& A_2^1\arrow[u,"d"']\arrow[l,"\partial"']\& A_3^1\arrow[u,"d"']\arrow[l,"\partial"']\&{}\arrow[l]\\
 0\& A_0\arrow[l]\& A_1^0\arrow[l,"\partial"']\arrow[u,"d"']\& A_2^0\arrow[l,"\partial"']\arrow[u,"d"']\& A_3^0\arrow[l,"\partial"']\arrow[u,"d"']\&\cdots\arrow[l]\\
 0\arrow[r]\& A_0\arrow[u,"\mathrm{id}","\varepsilon_0"']\& A_0\arrow[l,"0"']\arrow[u,"\varepsilon_1"']\& A_0\arrow[l,"\mathrm{id}"']\arrow[u,"\varepsilon_2"']\& A_0\arrow[l,"0"']\arrow[u,"\varepsilon_3"']\&\cdots\arrow[l,"\mathrm{id}"']
\end{tikzcd}}
% @22–25, revised set of conditions
\deleted{$\unsure{A_N^i=0|A_N^*=0}$, $H^N(A_N^\bullet)=0$ si $N\geq1$}
\[
 \left\{\begin{aligned}
 A_N^i&=0\quad\text{si }i\notin[0,N]\\
 H^N(A_N^*)&=0\quad\text{si }N\geq1\\
 H^{N-1}(A_N^*)&\ \deleted{\unsure{=}}\begin{cases}=0&\text{si }N\geq2\\
 \subset\operatorname{Ker}\delta&\text{si }N=1\end{cases}
 \end{aligned}\right.
\]
% @26–30, third remark
3) A cause de l'unicité, ou grâce à la version ``cat. abélienne'', on voit que les hom ``intégration'' $A_n^n\to A_0^0=A_0$ commutent à tous les endom des $A_*^*$, -- donc respectent une éventuelle structure de $k$-module
% archive sheet number: 55

\scanpage{57}
% typed verso
% @top centre, pagination
\mspage{- 155 -}
% @top right, running head
\hfill IX

% @01–05, graph and orientation
comme le conoyau de l'homomorphisme $\mathbb{Z}^{I\times J}\longrightarrow\mathbb{Z}^K$ induit par $h$, lequel est transposé de $\mathbb{Z}^{(h)}$. Notons d'ailleurs que le graphe $\Gamma(C)$ s'explicite aussi de façon évidente en termes de $h$, par $\Gamma_{\unsure{o|0}}=J$, $\Gamma_1=(K\times_I K-\Delta(K))/\text{symétrie}$. D'autre part, choisissant une orientation sur chacune des arrêtes de $\Gamma(C)$, et désignant par $\Gamma_1$ l'ensemble de ces arrêtes $(x_i,C',C'')$, on a
% @06, first homology
\[
 H_1(\Gamma(C),\mathbb{Z})=\operatorname{Ker}(d:\mathbb{Z}^{(\Gamma_1)}\longrightarrow\mathbb{Z}^{(J)})\tag*{(12.3.11)}
\]
% @07, differential
où l'opérateur différentiel est donné par
% @08, evaluation
\[
 d(x,C',C'')=g(C'')-g(C')\quad .\tag*{(12.3.12)}
\]
% @09–10, chain morphism
Ceci posé, l'homomorphisme (12.3.8) est induit par l'homomorphisme $v$ de complexes de chaînes :
% @11, complex map
\[
 v=(v_o,v_1):C_{\bullet}(\Gamma(C))=
 \big[\mathbb{Z}^{(\Gamma)}\longrightarrow\mathbb{Z}^{(J)}\big]\longrightarrow
 \big[\mathbb{Z}^{(K)}\longrightarrow\mathbb{Z}^{(I\times J)}\big]\tag*{(12.3.13)}
\]
% @12, components
donné par
% @13, values
\[
 v_1(x,C',C'')=(x,C'')-(x,C'),\quad v_o(\widetilde C_j)=(0,\widetilde C_j)\quad .\tag*{(12.3.14)}
\]
% @14–17, new subsection
12.4. Nous nous plaçons maintenant dans le cas envisagé à la fin de 12.2. En vertu de (12.1.12) et de (12.3.6) appliqué à $C=X_o$, le groupe des caractères $M$ du tore maximal $T_o$ de la fibre spéciale $A_o$ du modèle de Néron $A$ de $A=\underline{\mathrm{Pic}}^o_{X_\eta/\eta}$ est donné par
% @18, kernel
\[
 M=\operatorname{Ker}(\mathbb{Z}^I\xrightarrow{d}\mathbb{Z}^J)\quad ,\tag*{(12.4.1)}
\]
% @19–21, continuation
où $I$ est l'ensemble des points singuliers de $X_o$, et $J$ l'ensemble de ses composantes irréductibles, l'homomorphisme $\mathbb{Z}^I\longrightarrow\mathbb{Z}^J$ étant donné par
% @lower strip, pencil foliation
\margin{\handwritten{473}}
% @lower centre, printed label
\mspage{473}
% archive sheet number: 56

\scanpage{58}
% @01–02, category
Soit $(SS)$ la catégorie des ensembles semi-simpliciaux
% @03, presheaves
\[
 (SS)=\widehat{\mathbb{\Delta}}=\underline{\operatorname{Hom}}(\mathbb{\Delta}^{\circ},(\mathrm{Ens}))
\]
% @04, topos
C'est un topos, et si
% @05, object
\[
 \unsure{\mathscr{X}|\mathcal{X}|X}\in\operatorname{Ob}(SS)
\]
% @06–07, slice
est un ens $1/2$ simplicial, on regarde le topos induit
% @07, slice identification
\[
 (SS)/\mathscr{X}\simeq\widehat{\mathbb{\Delta}/\mathscr{X}}\simeq
 \underline{\operatorname{Hom}}(\mathbb{\Delta}^\circ,\mathrm{Ens})
\]
% @08–17, sheaves on the object
Pour nous, l'étude cohomologique de $\mathscr{X}$ est celle de ce topos induit. Les faisceaux \inserted{(d'ens)} ``sur $\mathscr{X}$'' sont les \deleted{\unsure{gr}} faisceaux semi-simpliciaux sur $\mathscr{F}\to\mathscr{X}$ sur $\mathscr{X}$, ou encore les foncteurs $(\mathbb{\Delta}/\mathscr{X})^\circ\to(\mathrm{Ens})$ ; les faisceaux abéliens sur $\mathscr{X}$ sont les objets groupes abéliens $F\to\mathscr{X}$ sur $\mathscr{X}$, ou encore les foncteurs $(\mathbb{\Delta}/\mathscr{X})^\circ\to(\mathrm{Ab})$ etc. Idem pour faisceaux abéliens gradués différentiels etc. Si $A_{\deleted{\unsure{*|\bullet}}}$ est un ens. ssl. on peut regarder
% @18, product
\[
 A\times\mathscr{X}\longrightarrow\mathscr{X}\quad\text{comme faisceau sur }\mathscr{X}
\]
% @19–21, induced sheaf
(c'est aussi le faisceau induit par $A_{\unsure{\alpha|*|\bullet}}$ sur le topos induit $(SS/\mathscr{X})$), c'est aussi le foncteur composé
% @22, composite
\[
 (\mathbb{\Delta}^{\circ}/\mathscr{X})^\circ\longrightarrow
 (\mathbb{\Delta})^\circ\xrightarrow{A}(\mathrm{Ens}).
\]
% archive sheet number: 57

\scanpage{59}
% typed verso
% @top centre, pagination
\mspage{- 154 -}
% @top right, running head
\hfill IX

% @01–08, graph of branches
de Lie, Chap IV, Annexe) $\Gamma(C)$ dont les sommets sont les composantes irréductibles de $\widetilde C$, deux sommets $\widetilde C_i$ et $\widetilde C_j$ étant reliés par un ensemble d'arrêtes en correspondance bijective avec l'ensemble $\widetilde C_i\times_C\widetilde C_j$ des ensembles de deux branches \uline{distinctes} de $C$ contenues dans $\widetilde C_i$ resp. $\widetilde C_j$ passant par le même point de $C$, ensemble égal à $\widetilde C_i\times_C\widetilde C_j$ si $i\ne j$, à $\widetilde C_i\times_C\widetilde C_i-\Delta(\widetilde C_i)$ modulo symétrie ($\Delta:\widetilde C\longrightarrow\widetilde C\times_C\widetilde C$ étant la diagonale) si $i=j$ (donc vide si $C_i$ est géométriquement unibranche). Ceci posé, on a des isomorphismes canoniques
% @09, canonical isomorphisms
\[
 M=D(T)\xleftarrow{\sim}H_1(\Gamma(C),\mathbb{Z})\quad,\quad
 \text{donc }\check M=R\simeq H^1(\Gamma(C),\mathbb{Z})\quad .\tag*{(12.3.7)}
\]
% @10–11, map
Nous n'aurons pas besoin de cette interprétation combinatoire de $M$, et allons nous borner à définir l'homomorphisme canonique
% @12, unnumbered display
\[
 H_1(\Gamma(C),\mathbb{Z})\longrightarrow M\quad ,
\]
% @13–18, sets of points and branches
laissant au lecteur le soin de prouver que c'est un isomorphisme. Soit donc $I$ l'ensemble des points de $C$ appartenant a au moins deux branches de $C$, $J$ (comme ci-dessus) l'ensemble des composantes irréductibles de $\widetilde C$ (ou de $C$, cela revient au même), $K$ l'ensemble des ``branches'' de $C$ en les points de $I$, i.e. l'image inverse de $I$ dans $\widetilde C$. On a donc une application canonique
% @19, application h
\[
 h:K\longrightarrow I\times J\quad,\quad h=(f,g)\quad ,\tag*{(12.3.9)}
\]
% @20, kernel introduction
et la description (12.3.4) de $M$ s'exprime en termes de $h$ comme
% @21, kernel display
\[
 M\simeq\operatorname{Ker}(\mathbb{Z}^{(h)}:\mathbb{Z}^{(K)}\longrightarrow\mathbb{Z}^{(I\times J)})\quad ,\tag*{(12.3.10)}
\]
% @22, continuation
car l'expression (12.3.3) de $R=\check M$ peut s'interprèter manifestement
% @lower strip, pencil foliation
\margin{\handwritten{472}}
% @lower centre, printed label
\mspage{472}
% archive sheet number: 58

\scanpage{60}
% @01–05, hypercohomology
Idem \deleted{pour} $A_*$ muni d'une structure suppl -- p ex. de groupe abélien, éventuellement diff et gradué. Dans ce dernier cas -- si les degrés sont bornés \unsure{infért} -- on peut donc prendre l'hypercohomologie.
% @06, definition
\[
 \mathbb{R}^*\unsure{\mathbb{\Gamma}|\Gamma}_{\unsure{?|,}}(\unsure{\mathscr{X}|\mathcal{X}|X},A_*^*)\overset{\mathrm{dfn}}{=}
 \mathbb{R}^*\Gamma_{\mathscr{X}}(\underline{A_{\mathscr{X}}^*})
\]
% @07–09, gloss below right term
\margin{complexe induit par $A_*^*$ sur $\mathscr{X}$ ou $A_*^*\times\mathscr{X}/\mathscr{X}$}
% @10–11, cochains
Soit $A$\inserted{$:(\deleted{\unsure{\Delta}}\mathbb{\Delta}/\mathscr{X})^{\circ}\longrightarrow(\mathrm{Ab})$} un groupe sur $\mathscr{X}$, considérons le complexe des cochaines sur $\mathscr{X}$ à coeff. ds $A$
% @12–13, product
\[
 C^n(\mathscr{X};A)\overset{\mathrm{dfn}}{=}
 \prod_{\substack{\sigma\in\operatorname{Ob}\,\Delta_n/\mathscr{X}\\=\mathscr{X}_n}}A(\sigma),
\]
% @14–15, final observation
Les $C^n$ forment un groupe co $1/2$ simplicial
% archive sheet number: 59

\scanpage{61}
% typed verso
% @top centre, pagination
\mspage{- 152 -}
% @top right, running head
\hfill IX

% @01–02, rank
$\underline{\mathrm{Pic}}^{\unsure{o|0}}_{\widetilde C/k}$ est de ``rang réductif'' nul (i.e. son tore maximal est nul), de sorte que le tore maximal de $\underline{\mathrm{Pic}}^o_{C/k}$ est contenu dans
% @03, kernel
\[
 N=\operatorname{Ker}(\underline{\mathrm{Pic}}_{C/k}\longrightarrow\underline{\mathrm{Pic}}_{\widetilde C/k})\quad ,
\]
% @04, continuation
donc est isomorphe canoniquement au tore maximal de ce dernier groupe.
% @05–14, description of the kernel
Or la structure de $N$ s'explicite aisément à l'aide de EGA IV 21.8.5 : on trouve que $N$ est canoniquement une extension d'un tore $T$ dont on va expliciter la description, par un groupe unipotent connexe $U$ (qui se dévisse de même, sur $k$, avec une suite de composition à quotients isomorphes au groupe additif $G_a$). D'ailleurs, si $T'$ est le tore maximal de $N$, alors $T'\cap U$ est à la fois unipotent et de type multiplicatif, donc nul (SGA 3 XVII 2.4), et l'image de $T'$ dans $T$ est identique à $T$ (SGA 3 XII 7.1 e)), de sorte que $T'\longrightarrow T$ est un isomorphisme. Donc $T$ est canoniquement isomorphe au tore maximal de $N$, donc de $\underline{\mathrm{Pic}}^o_{X/k}$.

% @15–21, branches
C'est pour expliciter la structure de $N$ qu'il sera commode de supposer $k$ séparablement clos, de sorte que les composantes irréductibles $C_i$ de $C$ sont géométriquement irréductibles, et que les points non normaux de $C$ sont radiciels sur $k$. Ceci dit, soit $J$ l'ensemble des composantes irréductibles de $\widetilde C$, et pour tout $x\in C$, soit $J(x)$ l'ensemble des ``branches'' de $C$ passant par $x$ (i.e. des points de $\widetilde C$ au-dessus de $x$) ; soient enfin
% @22, local quotient
\[
 R(x)=\mathbb{Z}^{J(x)}/\operatorname{Im}\mathbb{Z}\quad ,\tag*{(12.3.1)}
\]
% @23, diagonal
où $\mathbb{Z}\longrightarrow\mathbb{Z}^{J(x)}$ est l'application diagonale, et
% @24, quotient
\[
 R=(\bigoplus_x R_x)/\operatorname{Im}\mathbb{Z}^J\quad ,\tag*{(12.3.2)}
\]
% @lower strip, pencil foliation
\margin{\handwritten{470}}
% @lower centre, printed label
\mspage{470}
% archive sheet number: 60

\scanpage{62}
% @01, fibration
\[
 \Delta/\unsure{\mathscr X|\mathcal X}\xrightarrow{\text{fibration}}\Delta
\]
% @02–05, crossed diagram on the left
\diagram{\begin{tikzcd}[ampersand replacement=\&,row sep=small,column sep=small,execute at end picture={\draw[thin] (current bounding box.north east)--(current bounding box.south west);\draw[thin] (current bounding box.north west)--(current bounding box.south east);}]
\deleted{C}\arrow[rrr] \& \& \& \deleted{B}\\
\deleted{C^{\unsure{*|\circ}}}\arrow[d,"\deleted{F}"'] \& \deleted{C_b} \& \deleted{\unsure{i'(x')|i(x')}}\arrow[d] \& \deleted{b}\arrow[d,"\deleted{i}"]\\
\deleted{\deleted{\unsure{(\mathrm{Ab})}}} \& \deleted{C_{b'}}\arrow[u,"\deleted{\unsure{i^*|i^!}}"] \& \deleted{x'} \& \deleted{b'}\\
\deleted{\mathcal E} \& \& \&
\end{tikzcd}}
% @03–07, crossed diagrams on the right, upper comparison
\margin{\deleted{\unsure{$\Sigma$ direct|2 direct|si dual}}}
\diagram{\begin{tikzcd}[ampersand replacement=\&,row sep=large,column sep=small,execute at end picture={\draw[thin] (current bounding box.north west)--(current bounding box.south east);}]
\deleted{(b\in B)}\arrow[r,mapsto] \& \deleted{\deleted{\unsure{\prod}}\sum_{x\in C_b}F(x)}\arrow[r] \& \deleted{F(\unsure{i_!(x')|i^!(x')|i_*(x')})}\\
\deleted{(b'\in B)}\arrow[r,mapsto] \& \deleted{\sum_{x'\in C_{b'}}F(x')}\arrow[r]\arrow[u,"\deleted{\unsure{?|i}}"] \& \deleted{F(\deleted{\unsure{i}}x')}\arrow[u]
\end{tikzcd}}
% @08–10, crossed diagrams on the right, lower comparison
\diagram{\begin{tikzcd}[ampersand replacement=\&,row sep=large,column sep=small,execute at end picture={\draw[thin] (current bounding box.north east)--(current bounding box.south west);}]
\deleted{(b\in B)}\arrow[r,mapsto] \& \deleted{\prod_{x\in C_b}F(x)}\arrow[r]\arrow[d,"\deleted{\unsure{?}}"] \& \deleted{F(i^*(x'))}\arrow[d]\\
\deleted{(b'\in B)}\arrow[r,mapsto] \& \deleted{\prod_{x\in C_{b'}}F(x')}\arrow[r] \& \deleted{F(x')}
\end{tikzcd}}
% @09, diagonal near right of product square
\margin{\deleted{\unsure{OK|fibr}}}
% @11, left of next paragraph
\margin{\unsure{?}}
% @11–14, contravariant coefficients
Si $(\Delta/\mathscr X)^\circ\xrightarrow{A_*}\deleted{\unsure{(\mathrm{Ab})}}$\inserted{$\mathcal E$ (cat. avec sommes directes)} est un faisceau sur $\mathscr X$ (i.e. \deleted{faisceau} \inserted{syst.} de coeff. \uline{contravariant} pour un simplexe variable sur $\mathscr X$)
% @15–17, chains
on lui associe $C_*(\mathscr X,A_*)=i_{!}(A):(\Delta\deleted{\unsure{/\mathscr X}})^\circ\longrightarrow\mathcal E$ (où $i:\widehat{\Delta/\mathscr X}\longrightarrow\widehat{\Delta}$ est le morphisme d'induction des topos)
% @18–21, dual construction
Si $\Delta/\mathscr X\xrightarrow{A^*}\deleted{\unsure{(\mathrm{Ab})}}$\inserted{$\mathcal E$ (cat. avec produits)} est un ``cofaisceau'' sur $\mathscr X$ (i.e. syst. de coeff. \uuline{cov.} pour un simplexe variable de $\mathscr X$) on lui associe $C^*(\mathscr X,A^*)$ \deleted{\unsure{avec}} $:(\Delta\deleted{\unsure{/\mathscr X}})\longrightarrow\mathcal E$.
% @22–25, changing categories
On passe de l'un à l'autre en remplaçant $\mathcal E$ par $\mathcal E^{\circ}$ ; c'est compatible avec composition avec foncteur $\mathcal E\longrightarrow\mathcal E'$ qui commute aux sommes resp. produits.
% archive sheet number: 61

\scanpage{63}
% typed verso
% @top centre, pagination
\mspage{- 149 -}
% @top right, running head
\hfill IX

% @01–05, sections and divisors
\uline{en associant à tout} $i\in I$ \uline{la section} $S_i$ \uline{de} $P$ \uline{sur} $S$ (\uline{nulle en} $\eta$) \uline{définie par le faisceau inversible} $\underline O_X(X_{\unsure{o|0}}^i)$ \uline{sur} $S$. (NB \uline{Comme les anneaux locaux de} $X$ \uline{sont factoriels en hypothèse}, \uline{les cycles} $X_o^i$ \uline{sont bien des diviseurs}, \uline{ce qui donne un sens à} $\underline O_X(X_i^o)$ \uline{comme le faisceau inversible associé au diviseur en question} (EGA IV 21.2.8.1).) \uline{On a}
% @06, intersection
\[
 P^o\cap E=\text{\uline{section nulle de} }P\quad .\tag*{(12.1.6)}
\]
% @07, quotient
c) \uline{Sous les conditions de} b), \uline{le quotient}
% @08, quotient formula
\[
 Q=P/E\tag*{(12.1.7)}
\]
% @09–11, representability
\uline{est représentable par un schéma en groupes lisse et séparé sur} $S$, \uline{et ce dernier est} néronien, \uline{i}.\uline{e}. \uline{satisfait la propriété universelle} (1.1.2) \uline{en termes de sa fibre spéciale}.

% @12–14, exact sequence
d) \uline{Sous les conditions de} c), \uline{supposons de plus que la fibre géométrique générique de} $X$ \uline{soit} intègre. \uline{Alors on a une suite exacte canonique}
% @15, degree
\[
 0\longrightarrow A\longrightarrow Q\xrightarrow{\deg}\mathbb{Z}_S\longrightarrow 0\quad ,\tag*{(12.1.8)}
\]
% @16–18, Néron model
\uline{où} $\deg$ \uline{est induit par l'homomorphisme} ``\uline{degré}'' $P\longrightarrow\mathbb{Z}_S$, \uline{et où} $A$ \uline{est un schéma en groupes lisse séparé et} de type fini \uline{sur} $S$, \uline{qui est également néronien} \uline{i}.\uline{e}. \uline{qui est un modèle de Néron de}
% @19, generic fibre
\[
 A_\eta=P_\eta^o=\underline{\mathrm{Pic}}_{X_\eta/\eta}^{o}\quad .\tag*{(12.1.9)}
\]
% @20, composite
\uline{D'autre part}, \uline{l'homomorphisme composé} $P^o\longrightarrow P\to Q$ \uline{induit un isomorphisme}
% @21, identity components
\[
 P^o\xrightarrow{\sim}A^o\quad .\tag*{(12.1.10)}
\]
% @lower strip, pencil foliation
\margin{\handwritten{467}}
% @lower centre, printed label
\mspage{467}
% archive sheet number: 62

